\section{Introduction}

Model-Based Systems Engineering (MBSE) promises improved traceability and earlier verification, yet many requirements remain in natural language or diagrammatic forms \parencite{salado2019mbse}. Checking whether an implementation satisfies stated requirements depends on manual reconstruction of intent: systems engineers cannot mechanically determine conformance without informal argument. Encoding requirements as temporal properties over explicit system trajectories is one path toward model-based requirements that are formally analyzable and executable rather than purely documentation \parencite{salado2019mbse}.

Wymore's Tricotyledon Theory of System Design (T3SD) offers a rigorous foundation to build formal requirement: cotyledon-structured design spaces, trajectory theorems, and constructive verification via reference systems and homomorphisms \parencite{Wymore1993,salado2019wymore}. As generic model checking lacks T3SD's design-space semantics, we reinterpret the Functionality Cotyledon (FC) assertionally while preserving cotyledon structure. Constructing and maintaining homomorphisms is largely manual and poorly matched to evolving implementations, especially when concurrent subsystems induce large proof obligations \parencite{Wymore1993}. In our experiments, even modest system homomorphism arguments span thousands of lines of interactive proof script across conjunctive composition and dynamic reconfiguration; Section~\ref{sec:automated-checking} contrasts this cost with automated solver-based checking of the same boundary properties.

We propose encoding functional requirements as a property set $\Phi$ in temporal logic dialect $\mathcal{L}$ and verifying a candidate implementation by checking $Z_{\text{impl}} \models_{\mathcal{L}} \Phi$ over admissible Wymore trajectories. Furthermore, any Wymore DSYSTEM $Z$ can be translated into an FO-LTL formula $\psi_Z$ that encodes its execution semantics so requirement verification can reuse the same vocabulary as $\Phi$ and, in suitable fragments, automated model checkers or SMT solvers \parencite{BaierKatoen2008,deMouraZ32008}. Importantly, assertional FC does not require a concrete $Z_{\text{spec}}$ for checking.

\miniheading{Contributions}
\begin{itemize}
    \item A parameterized assertional reinterpretation of the Functionality Cotyledon, reframing functional verification as temporal property checking over Wymore trajectories rather than as construction of Wymore system homomorphisms.
    \item Trace-based verification semantics for state and output trajectories, aligned with Wymore's trajectory theorems \parencite{Wymore1993}, supporting model-based requirements executable against open-system boundaries \parencite{salado2019mbse}.
    \item A canonical FO-LTL encoding of Wymore DSYSTEM execution semantics.
    \item A workflow comparison of constructive and assertional verification methods (case study in Section~\ref{sec:case-study}).
\end{itemize}

The paper proceeds as follows: Section \ref{sec:background} summarizes T3SD, notation, and temporal logics; Section \ref{sec:methodology} defines trace semantics, the DSYSTEM-to-FO-LTL encoding, and verification workflows; Section \ref{sec:case-study} compares constructive and assertional workflows on an illustrative system; Section \ref{sec:discussion} relates assertional FC to classical T3SD. Appendix~\ref{app:hom-biimp} walks through the headline dynamics-encoding bi-implication (hom projection); Appendix~\ref{app:encoding} proves the execution-encoding theorem for FO-LTL.
